%! Characteristics of Linear Functions — Unit 1, Lesson 1.0 — Student Packet Cover Sheet
\documentclass[10pt]{article}
\usepackage{algebra2-article}
\usepackage{algebra2-boxes}

\begin{document}

% Full-bleed forest banner
\begin{tikzpicture}[remember picture, overlay]
  \fill[forest] (current page.north west)
    rectangle ([yshift=-0.9in]current page.north east);
\end{tikzpicture}
\vspace{-0.8in}

\begin{center}
  {\color{white}\textbf{\LARGE Algebra 2}}\\[4pt]
  {\color{forestmid}\large\textbf{Unit 1 \quad Linear Functions}}\\[6pt]
  {\color{white}\textbf{\large Lesson 1.0 \quad Characteristics of Linear Functions}}
\end{center}

\vspace{0.22in}
\namedateperiod
\vspace{0.18in}

% Targets are deliberately spoiler-free (EFFL): they describe what students will be able to
% DO in plain language, without pre-naming the vocabulary the debrief will attach.
\begin{learningtargetbox}
\small
\begin{enumerate}[leftmargin=1.5em, itemsep=5pt]
  \item \ldots{} read the story a straight-line graph tells --- where it starts, where it hits
        zero, and how fast it changes --- and say what each of those means in the situation.
  \item \ldots{} tell the difference between \emph{where a graph sits} and \emph{which way it is
        heading} --- and explain why those are two different questions.
  \item \ldots{} decide which inputs and outputs make sense --- for the math on its own, and for
        the story it models.
\end{enumerate}
\end{learningtargetbox}

\vspace{0.14in}

\begin{tocbox}
\small
\renewcommand{\arraystretch}{1.5}
\begin{tabularx}{\linewidth}{@{} c l X r @{}}
\rowcolor{forestbg}
\textbf{\#} & \textbf{Component} & \textbf{Description} & \textbf{Score} \\
\midrule
1 & Warm-Up    & Getting ready: reading points, evaluating $g(x)$, spotting a steady pattern & \blank{1.2cm} \\
2 & Experience \& Formalize & \emph{Freezing Point} --- two temperature stories (with your
                 group), then \mbox{QuickNotes} and an application we work together          & \blank{1.2cm} \\
3 & Check Your Understanding & Six practice exercises --- \emph{this is the lesson's practice,
                 so it is not scored}                                                         & \textbf{NA} \\
4 & Homework   & In \textbf{DeltaMath} --- \emph{due Sunday, Aug.\ 30, 11:59\,pm}              & \blank{1.2cm} \\
\midrule
  & \multicolumn{2}{r}{\textbf{Total}} & \blank{1.2cm} \\
\end{tabularx}
\end{tocbox}

\vspace{0.10in}

\begin{remindbox}
\small In this lesson you will \textbf{experience first, formalize later}: you and your group will
work out two temperature stories using what you already know --- tables, rules, and points on a
graph --- and only afterward will we name what you found. The Check Your Understanding exercises
are \textbf{practice, not a graded quiz} --- your graded practice is the DeltaMath set.
\end{remindbox}

\end{document}
